\scanpage{62}
\uline{Corollaire} 2.\unsure{?|2}. \margin{\rule[-2.8\baselineskip]{0.45pt}{3.5\baselineskip}}Notations de 2.1. Considérons $ {}^t(u_Z):B_Y\longrightarrow A_X$, alors $ {}^t(u_Z)=u_{({}^tZ)}$, i.e. on a pour tout diviseur $D$ alg. équiv. à $0$ sur $Y$, tel que $Z(D)$ soit défini,
\[
S_X(Z(D))={}^t(u_Z)(D)
\]
% slanted
\margin{Attention aux signes éventuels !}
Vérification par exemple par voie cohomologique.

$[$ \uline{Proposition} 2.3. Conditions de 1.1. \deleted{Alors on a}\inserted{Considérons} $v_Z={}^t(u_Z):B_X\longrightarrow B_Y$. Alors on a, pour tout diviseur alg. équiv. à $0$ $D$ sur $X$, tel que $Z(D)$ soit défini,
% slanted
\margin{devrait remonter au N\textsuperscript{o}1.}
\[
{}^tZ(D)={}^t(\unsure{u|v}_Z)(D)
\]

N.B. Ceci doit pouvoir servir à une autre démonstration de 1.1., via un hom de schémas en groupes $\underline{\operatorname{Pic}}_{X/k}\longrightarrow\underline{\operatorname{Pic}}_{Y/k}$ !
% slanted
\margin{attention aux signes éventuels !}

Vérification de 2.3. par AQT $]$

\uline{Proposition} 2.4. Avec les notations de 2.1. \deleted{Supposons que}\inserted{Fixons $l$,} pour \deleted{un $l$ convenable,}\inserted{que} l'hom. $a(\gamma(z)):\allowbreak H^1(Y,\mathbb Q_l)\longrightarrow H^{2n-1}(X,\mathbb Q_l)(n-1)$ soit un isom \deleted{\unsure{modulo|mod.} torsion}\inserted{il faut et s. que $u_Z$ soit une isogénie (cond. ind. de $l$)} \deleted{\unsure{diviseurs|et suffisante}}. Alors il existe un \deleted{cycle}diviseur
% archive sheet number: 61

